% ============================================================================
% LANGENTWURF LE-017: Fotos & Dateien
% Profil: S (Senioren 65+) | Tier: 2 (Standard)
% ============================================================================

\documentclass[11pt,a4paper]{article}
\usepackage[utf8]{inputenc}
\usepackage[T1]{fontenc}
\usepackage[ngerman]{babel}
\usepackage{geometry}
\usepackage{fancyhdr}
\usepackage{lastpage}
\usepackage{xcolor}
\usepackage{tcolorbox}
\usepackage{enumitem}
\usepackage{longtable}
\usepackage{hyperref}
\usepackage{amssymb}

\geometry{left=2.5cm, right=2.5cm, top=2.5cm, bottom=2.5cm}
\definecolor{fach}{HTML}{b35c00}
\definecolor{gray}{HTML}{666666}
\definecolor{successgreen}{HTML}{2a7a4b}

\tcbuselibrary{skins,breakable}
\newtcolorbox{scorebox}{ colback=white, colframe=fach, fonttitle=\bfseries, title=Didaktik-Score, breakable }

\pagestyle{fancy}
\fancyhf{}
\fancyhead[L]{\small\textcolor{gray}{Digitale Grundlagen -- 017 Fotos \& Dateien}}
\fancyhead[R]{\small\textcolor{gray}{LE Tier 2 / Profil S}}
\fancyfoot[C]{\small\thepage{} / \pageref{LastPage}}
\renewcommand{\headrulewidth}{0.4pt}

\begin{document}

\begin{titlepage}
    \centering
    \vspace*{2cm}
    {\Large\textcolor{gray}{Digitale Grundlagen -- Senioren 65+}}\\[2cm]
    {\Huge\textcolor{fach}{\textbf{Langentwurf}}}\\[0.5cm]
    {\large Tier 2 | Profil S}\\[2cm]
    {\LARGE\textbf{017 -- Fotos \& Dateien verwalten}}\\[0.5cm]
    {\large Modul D: Produktivität}\\[2cm]
    \begin{tabular}{rl}
        \textbf{Zeitrahmen:} & 75--90 Minuten \\
        \textbf{Vorkenntnisse:} & Kamera nutzen, Notizen \\
    \end{tabular}
    \vfill
    {\normalsize Stand: \today}
\end{titlepage}

\tableofcontents
\newpage

\section{Bedingungsanalyse}

\subsection{Ausgangssituation}
\begin{itemize}
    \item Tausende Fotos, keine Ordnung
    \item ``Speicher voll'' -- keine Ahnung warum
    \item Fotos finden ist Glückssache
    \item Unterschied Fotos-App vs. Dateien-App unklar
    \item Angst vor Löschen (``Ist es dann weg?'')
\end{itemize}

\subsection{Typische Probleme}
\begin{itemize}
    \item ``Wo ist das Foto vom letzten Urlaub?''
    \item ``Mein Speicher ist voll''
    \item ``Wie lösche ich Fotos, ohne sie ganz zu verlieren?''
    \item ``Was ist die Dateien-App?''
\end{itemize}

\section{Sachanalyse}

\subsection{Kernkonzepte}
\begin{enumerate}
    \item \textbf{Fotos-App:} Automatische Sammlung aller Bilder und Videos
    \item \textbf{Alben:} Fotos in Gruppen sortieren (z.B. ``Familie'', ``Urlaub'')
    \item \textbf{Suche:} Fotos nach Ort, Datum, Inhalt finden
    \item \textbf{Dateien-App:} Für Dokumente, PDFs, Downloads
    \item \textbf{Zuletzt gelöscht:} 30 Tage Rettungschance
\end{enumerate}

\subsection{Alltagsanalogien}
\begin{tabular}{|l|p{8cm}|}
\hline
\textbf{Begriff} & \textbf{Analogie} \\
\hline
Fotos-App & Fotoalbum, das sich selbst füllt \\
Alben & Verschiedene Alben im Regal (Familie, Urlaub...) \\
Suche & Bibliothekar, der Fotos nach Beschreibung findet \\
Dateien-App & Aktenschrank für Dokumente \\
Zuletzt gelöscht & Papierkorb -- Inhalt noch 30 Tage drin \\
\hline
\end{tabular}

\section{Didaktische Analyse}

\subsection{Stellung in der Reihe}
Wichtig für Speichermanagement und Organisation. Vorbereitung auf Cloud (018).

\subsection{Didaktische Reduktion}
\textbf{Ausgeklammert:} Fotobearbeitung, Live Photos, Speicheroptimierung im Detail.

\textbf{Fokus:} Fotos finden, Alben erstellen, Platz schaffen.

\section{Lernziele}

\begin{tabular}{|c|p{7cm}|c|c|}
\hline
\textbf{Nr.} & \textbf{Die Teilnehmer können...} & \textbf{Stufe} & \textbf{Niveau} \\
\hline
FZ1 & die Fotos-App navigieren (Mediathek, Alben) & Anwenden & Basis \\
FZ2 & ein Foto über die Suche finden & Anwenden & Basis \\
FZ3 & ein Album erstellen und Fotos hinzufügen & Anwenden & Basis \\
FZ4 & Fotos löschen und aus ``Zuletzt gelöscht'' wiederherstellen & Anwenden & Basis \\
FZ5 & die Dateien-App für Dokumente nutzen & Anwenden & Vertiefung \\
\hline
\end{tabular}

\section{Verlaufsplanung}

\begin{longtable}{|p{1cm}|p{2cm}|p{5cm}|p{3cm}|p{2cm}|}
\hline
\textbf{Zeit} & \textbf{Phase} & \textbf{Inhalt} & \textbf{TN-Aktivität} & \textbf{Material} \\
\hline
0--10 & Einstieg & ``Wie viele Fotos haben Sie?'' Problem aufzeigen & Schätzen & -- \\
\hline
10--25 & Fotos-App & Mediathek, Für dich, Alben erklären, Navigation & Mitmachen & S.1 \\
\hline
25--40 & Suche & Nach Ort, Datum, ``Hund'', ``Strand'' suchen & Selbst suchen & S.1 \\
\hline
40--55 & Alben & Album erstellen, Fotos auswählen und hinzufügen & Üben & S.2 \\
\hline
55--70 & Löschen & Foto löschen, Zuletzt gelöscht, Wiederherstellen & Ausprobieren & S.2 \\
\hline
70--85 & Dateien & Dateien-App kurz vorstellen, Downloads finden & Anschauen & S.3 \\
\hline
\end{longtable}

\section{Lösungen}

\textbf{Foto suchen:}
\begin{itemize}
    \item Fotos-App öffnen
    \item Unten auf ``Suchen'' tippen
    \item Stichwort eingeben: Ort (``Berlin''), Zeit (``2023''), Inhalt (``Katze'')
    \item iPhone erkennt automatisch, was auf Fotos ist!
\end{itemize}

\textbf{Album erstellen:}
\begin{itemize}
    \item Alben $\rightarrow$ + oben links $\rightarrow$ ``Neues Album''
    \item Namen eingeben (z.B. ``Enkel'')
    \item Fotos auswählen $\rightarrow$ Hinzufügen
    \item Alternative: Foto antippen $\rightarrow$ Teilen $\rightarrow$ Zum Album hinzufügen
\end{itemize}

\textbf{Fotos löschen (sicher):}
\begin{itemize}
    \item Foto(s) auswählen
    \item Papierkorb-Symbol antippen
    \item Foto landet in ``Zuletzt gelöscht'' (30 Tage)
    \item Dort: Wiederherstellen ODER endgültig löschen
\end{itemize}

\textbf{Speicherplatz-Tipps:}
\begin{itemize}
    \item Doppelte Fotos löschen (unter Alben $\rightarrow$ Duplikate)
    \item Große Videos prüfen
    \item ``Zuletzt gelöscht'' leeren
    \item iCloud-Fotos aktivieren (siehe 018)
\end{itemize}

\section{Didaktik-Score}

\begin{scorebox}
\begin{tabular}{|c|l|c|c|p{3.5cm}|}
\hline
\# & Dimension & Gew. & Score & Notiz \\
\hline
1 & Differenzierung & 15\% & 8/10 & Basis gut, Vertiefung ok \\
2 & Taxonomie & 10\% & 9/10 & Anwenden stark \\
3 & Alltagsrelevanz & 10\% & 10/10 & Jeder hat Fotos! \\
4 & Aktivierung & 10\% & 10/10 & Eigene Fotos nutzen \\
5 & Scaffolding & 10\% & 9/10 & Schritt-für-Schritt \\
6 & Feedback & 10\% & 8/10 & Checkliste \\
7 & Sprachliche Klarheit & 10\% & 9/10 & Gute Analogien \\
8 & Motivation & 10\% & 10/10 & Ordnung motiviert \\
9 & Barrierefreiheit & 10\% & 9/10 & Große Darstellung \\
10 & Praktikabilität & 5\% & 8/10 & 85 Min evtl. lang \\
\hline
\multicolumn{3}{|r|}{\textbf{GESAMT:}} & \textbf{9.15/10} & \textcolor{successgreen}{\textbf{FREIGABE}} \\
\hline
\end{tabular}
\end{scorebox}

\end{document}
